% 1_resumen.tex

\begin{abstract}
Este corresponde al primer avance del proyecto del curso IE0527: Ingeniería de Comunicaciones, el cual consiste en transmisión inalámbrica de audio.
En el presente, se incluye el análisis del problema a resolver y la investigación teórica que fundamenta el diseño del sistema.
Se identificaron los principales desafíos tecnológicos del sistema, entre los que destacan la fragmentación del audio en paquetes de tamaño restringido, la sincronización autónoma del receptor, el manejo de ruido y la confiabilidad del enlace inalámbrico bajo restricciones de latencia impuestas por el enunciado.

Asimismo, se investigaron los principios de digitalización de audio, incluyendo frecuencia de muestreo, cuantización, formatos digitales y esquemas de codificación, con énfasis en su impacto sobre el volumen de datos y la viabilidad de la transmisión.
Adicionalmente, se analizaron las limitaciones del canal inalámbrico en la banda ISM de 2.4 GHz, considerando fenómenos como atenuación, interferencia, \textit{shadowing} y capacidad máxima del canal según el teorema de Shannon.
Los resultados de esta investigación establecen las bases teóricas para la selección del hardware y el diseño del protocolo de transmisión, que se desarrollarán en el siguiente avance.
\end{abstract}

\begin{IEEEkeywords}
% En orden alfabético
Ancho de banda, audio digital, banda ISM, codificación de audio, digitalización, frecuencia de muestreo, sincronización, transmisión inalámbrica.
\end{IEEEkeywords}